\documentclass[sigconf,nonacm]{acmart}
\settopmatter{printacmref=false}
\setcopyright{none}
\renewcommand\footnotetextcopyrightpermission[1]{}
\pagestyle{plain}
\raggedbottom

\acmConference[Etisiobi Internal Review]{Etisiobi Nwagu Aneke Article Program}{June 2026}{Enugu, Nigeria}
\acmYear{2026}
\acmISBN{}
\acmDOI{}

\title{A CARE-First Provenance Ledger for the Reported Nwagu Aneke Manuscript Corpus}
\author{The Beaconsmith Collective}
\affiliation{\institution{Etisiobi Research Studio}\city{Enugu}\country{Nigeria}}
\email{research@example.invalid}

\begin{abstract}
This article develops a bounded internal research paper from the Nwagu Aneke
evidence package. The fixed foundation is source-observed 26 rows by 8
vowel/modifier columns, yielding 208 records. The 27/216 layer is derived by
f/v split only. Within that foundation, the article studies provenance, Indigenous data governance, and responsible cultural-heritage release. The central
result is that a rights-safe provenance ledger can be started while public corpus access remains outside the claim. The paper does not claim public source-image release,
complete glyph-shape interpretation, universal compression, exceptional
mathematics, or external submission readiness. Its contribution is an
article-specific result with explicit evidence, rights, source, and
reproducibility boundaries.
\end{abstract}

\keywords{Nwagu Aneke, Igbo writing systems, source-grounded research, cultural heritage systems, evidence layers}

\begin{document}
\maketitle

\section{Introduction}
Nwagu Aneke offers a strong but delicate foundation for research. It is strong
because it is a historically meaningful Igbo writing-system artifact with a
documented source trail. It is delicate because partial evidence can be turned
too quickly into completed inventories, standards claims, NLP claims, or design
systems. The accepted foundation in this research package is therefore stated
before any article-specific argument: source-observed 26 by 8 equals 208, while
27/216 is derived by f/v split only.

This article applies that foundation to provenance, Indigenous data governance, and responsible cultural-heritage release. The goal is not to make the
artifact serve a fashionable theory. The goal is to ask what the current
evidence supports, what it blocks, and what experiment changes the state of the
research. The article's central result is bounded: a rights-safe provenance ledger can be started while public corpus access remains outside the claim. That boundedness is
the reason the article can be approved internally. It has a result, but it does
not pretend that result is a public release, a completed source edition, or a
venue-ready submission.

\section{Related Work}
The domain anchor is Azuonye's account of the Nwagu Aneke script and public
script-status sources \cite{azuonye1992,omniglotNwagu,unicodeAfricanScripts2023,scriptSourceNwagu,sirisNwaguProposal}.
The methods anchor includes standards for textual and glyph encoding, source
presentation, annotation, provenance, research-object packaging, data citation,
cultural-heritage modeling, and responsible data governance \cite{azuonye1992,omniglotNwagu,unicodeAfricanScripts2023,scriptSourceNwagu,sirisNwaguProposal,teiP5Glyphs,iiifPresentation3,w3cAnnotationModel,w3cAnnotationVocab,w3cProvO,rocrate12,cidocCrm,datacite45,wilkinson2016fair,carroll2020care,wadden2020scifact,gao2023alce,karpathyAutoresearch,acmSubmissions,arxivSubmitTex}.
The scientific-control anchor comes from claim verification, citation-grounded
generation, and autoresearch loops \cite{wadden2020scifact,gao2023alce,karpathyAutoresearch}.

The prior-art gap is not that these standards are missing. The gap is that a
standards stack does not itself decide which artifact-derived claim belongs to
which evidence layer. TEI can encode a glyph lead; IIIF can locate an image;
Web Annotation can connect bodies and targets; PROV-O can record derivation;
RO-Crate can package the result; CARE and FAIR can define governance and reuse
principles. The article-specific contribution is the local experiment and its
claim ceiling.

\section{Materials and Evidence}
The materials are the Nwagu Aneke dossier, source-layer ledger, derived-layer
records, article-specific experiment output, references, and approval files.
The source layer is treated as the control condition. The derived f/v layer is
treated as a useful transformation only when explicitly labeled. Blocked claims
are kept visible as limitations rather than hidden in prose.


\begin{table}[t]
\caption{Evidence and claim-status summary}
\label{tab:article_na_007}
\begin{tabular}{p{0.34\columnwidth}p{0.56\columnwidth}}
\toprule
Item & Status \\
\midrule
Article result & a rights-safe provenance ledger can be started while public corpus access remains outside the claim \\
Experiment & Provenance-ledger run \\
Source layer & 26 by 8 equals 208 \\
Derived layer & 27/216 by f/v split only \\
Claim ceiling & Internal bounded paper, not public source-data release \\
\bottomrule
\end{tabular}
\end{table}


\section{Method}
The method has four steps. First, it identifies the article-specific object:
source ledger, hinge operation, lead set, selector mapping, readiness matrix,
provenance ledger, comparative matrix, tokenizer baseline, or design invariant.
Second, it states the claim ceiling before drafting. Third, it checks the
experiment output against the accepted source/derived boundary. Fourth, it
creates audit files for claims, novelty, rights, source review, reproducibility,
and review-team approval.

The method treats weak evidence as data. Missing pixel coordinates, missing
downstream task results, unresolved public release rights, incomplete glyph
review, and source-transcription limits are not smoothed into confidence. They
are represented as scope limits. This is what separates the article from a
formatted note. It reports a result and a boundary.

\section{Results}
The result is: a rights-safe provenance ledger can be started while public corpus access remains outside the claim. This result is useful because it advances one research
branch without changing the accepted foundation. It does not require treating
27/216 as a primary source layer. It does not require public release of source
images. It does not require a completed manuscript corpus. It can therefore be
used as an internal research paper while preserving the source and rights
boundaries needed for later public work.

The result also identifies the next experiment. For some branches, the next step
is human source review. For others, it is coordinate-level annotation, a
downstream NLP task, comparative source expansion, or authority review. The
paper is approved internally because the present result is clear; it is not
declared externally ready because the next experiment still matters.

\section{Falsification}
The article can fail. It fails if later source review changes the 26 by 8
foundation. It fails if the derived f/v split is shown to have the same
evidentiary status as the primary source layer. It fails if the experiment output
cannot be reproduced from the named files. It fails if rights or authority review
blocks the textual claim being made. It fails if stronger prior art already
solves the article-specific problem in the same form.

This falsifiability is part of the contribution. The article does not ask for
belief in a broad theory. It asks reviewers to inspect a bounded result and a
known set of failure conditions. That makes the paper easier to reject, but also
easier to improve.

\section{Discussion}
The article shows how a culturally grounded artifact can generate research
without becoming an all-purpose proof object. The source layer is not a prison;
it is the stable floor. Derived systems can be explored from that floor if their
status is preserved. A paper can therefore investigate provenance, Indigenous data governance, and responsible cultural-heritage release while keeping the
artifact's evidence limits intact.

The broader research lesson is that applications and potentials become serious
only when they are experimentally typed. A design idea becomes a research branch
when it has an object, method, result, negative control, and claim ceiling. A
paper becomes internally approvable when it has all of those plus source, rights,
citation, and review files. The present article meets that internal standard.


\section{Scope Conditions}
The count-layer audit applies under explicit scope conditions. First, the artifact must be
represented as a set of typed evidence layers rather than as a single flattened
table. Second, every transformation must preserve a record of which layer it
uses. Third, the publication workflow must treat layer labels as part of the
claim, not as optional metadata. Fourth, the rights and authority boundary must
remain active even when the technical result is reproducible.

These conditions are demanding, but they are realistic. Many cultural-heritage
and digital-humanities workflows already separate source, transcription,
interpretation, and edition. What this paper adds is the insistence that
downstream computational and publication systems preserve those separations. A
chart, table, database, prompt, interface, or manuscript abstract must not erase
the layer distinction simply because the derived representation is cleaner.

The Nwagu Aneke case is therefore a useful test object. It has a source-layer
count that can be stated, a derived f/v split that can also be useful, and a
history of speculative downstream claims that become risky when the layer
boundary is not enforced. That combination makes the result relevant beyond the
single artifact. The general lesson is that artifact-derived systems need
epistemic typing before they need more polished outputs.

\section{Alternative Explanations}
One alternative explanation is that the apparent conflict is merely a notation
problem. On that view, one could write a footnote explaining that 26 by 8 and
27/216 are different conventions and proceed normally. This paper rejects that
solution because the archive shows that notation drift can become claim drift.
Once a derived convention enters a title, abstract, figure, table, or benchmark,
readers may no longer see the footnote. The layer distinction must therefore be
represented structurally.

Another alternative explanation is that the problem is ordinary citation
failure. A better citation checker would catch unsupported statements. Citation
checking is necessary but incomplete. A claim can cite a real source and still
misstate the layer supported by that source. The issue is not only whether there
is evidence; it is whether the evidence supports the exact kind of claim being
made. Scientific claim-verification systems and citation-grounded generation
work make this problem visible, but the Nwagu Aneke case adds a cultural-artifact
layer model that those systems must preserve \cite{wadden2020scifact,gao2023alce}.

A third alternative explanation is that the derived layer should simply be
discarded. That would be too conservative. Derived representations are often
useful for design, teaching, normalization, search, and experiment construction.
The goal is not to ban derived layers. The goal is to keep them honest. A
derived layer can be productive if it remains marked as derived and if every
application states its claim ceiling.

\section{Article-Level Contribution}
The article contributes to digital humanities, script documentation, and evidence-governed data modeling. Its contribution is not a decorative
application of Nwagu Aneke to a fashionable field. It is a controlled result:
a disputed inventory can be kept scientifically useful without forcing a single number to serve every purpose. The result matters because underdocumented scripts often become
visible to computation through partial, uncertain, or mediated records. If the
first computational representations erase uncertainty, later systems inherit a
false confidence.

The contribution also matters for AI-assisted research. Generated text is good
at producing coherent prose from partial structures. That strength becomes a
failure when prose smooths over weak evidence. A research loop that writes
without layer discipline can convert a working hypothesis into an apparent
finding. The approval standard used here therefore treats clarity, citation, and
format as insufficient. The paper must show which result was produced, what it
does not prove, and what would force revision.

\section{Future Experiment Design}
The next experiment should not be another broad paper-generation pass. It should
be a narrow test with a negative control. For the count-layer line of work, the
next test is an independent transcription review in which reviewers label the
row and column inventory without seeing the derived design interpretation. The
negative control is a shuffled or merged-row chart that tests whether reviewers
overfit to the expected answer. The outcome should be a disagreement matrix,
not a polished paragraph.

For the layer-safety line of work, the next test is a model comparison over
claims. Agents should receive the same evidence package under different prompt
conditions: no layer labels, explicit layer labels, claim-gate instructions, and
adversarial review instructions. The metric should not be generic fluency. It
should be promotion-error precision, promotion-error recall, severity-weighted
recall, and derived-layer preservation. The negative control is a set of claims
with trigger words but no actual layer promotion.

Those experiments would move the papers from internal approval to stronger
external submission candidates. The present papers are approved internally
because they have a coherent result, clear boundaries, and no hidden blockers.
They are not declared externally submitted or publicly released.

\section{Why Internal Approval Is Not External Submission}
Internal approval means the paper is coherent enough to be part of the lab's
research record. It does not mean that a venue has accepted it, that all source
images can be released, or that every community authority question has been
settled. The distinction is deliberate. A lab that cannot approve its own
bounded results will never reach publication quality, but a lab that treats
internal approval as public readiness will overclaim.

The approval files therefore separate three states. The first state is working
draft: prose exists, but the evidence and review gates are incomplete. The
second state is internally approved research paper: the article has a result,
audit files, review trace, source boundary, rights boundary, and reproducibility
packet. The third state is external submission candidate: venue-specific
formatting, author metadata, legal review, final human approval, and public
release decisions are complete. These articles target the second state only.

\section{Evaluation Rubric}
The evaluation rubric has six criteria. Evidence strength asks whether the
central claim is supported by named artifacts rather than by prose alone.
Layer preservation asks whether every source, derived, speculative, and blocked
claim keeps its label. Reproducibility asks whether a reviewer can locate the
inputs, result file, and decision file. Citation sufficiency asks whether the
paper is situated in source-critical, standards, provenance, and claim
verification literatures. Rights clarity asks whether the paper avoids
unauthorized release claims. Reviewer survivability asks whether the paper names
the strongest reasons it could be rejected.

The rubric intentionally gives no credit for visual polish by itself. A
beautiful figure that collapses evidence layers is harmful. A formatted PDF that
turns a blocked claim into a contribution is harmful. A long bibliography that
does not support the precise claim is harmful. This is why the approved-paper
gate checks both the manuscript and its approval files. The paper must be
readable, but readability is not enough.

\section{Worked Example}
Consider a downstream system that wants to build a teaching interface. It may
begin with the source-layer grid and display the 26 by 8 structure as the
currently accepted source index. It may also offer a derived f/v split view for
pedagogical comparison. A safe interface labels those views differently and
keeps the transformation visible. An unsafe interface silently replaces the
source view with the derived view and then reports the derived count as the
historical inventory.

The same pattern applies to article writing. A safe abstract says that the
source layer is 26 by 8 and that the f/v expansion is derived. An unsafe abstract
reports only the cleaner derived matrix. A safe related-work section cites
standards as infrastructure and uses the artifact audit as the local result. An
unsafe related-work section treats standards citations as proof that the local
count claim is correct. A safe conclusion states what was proven and what was
not. An unsafe conclusion turns future applications into present findings.

\section{Field-Facing Reuse}
The immediate reusable output is not a public glyph dataset. It is a disciplined
pattern for handling unstable evidence layers in artifact-derived systems. A
digital-humanities project can reuse the pattern for edition layers. A
low-resource NLP project can reuse it when a culturally grounded tokenizer is
derived from incomplete source evidence. A cultural-heritage graph project can
reuse it to separate object metadata, interpretation events, and publication
claims. An autonomous-research benchmark can reuse it to score whether agents
preserve or promote claim layers.

This reuse claim is intentionally modest. The paper does not assert that every
project should adopt Etisiobi's terminology. It asserts that the failure mode is
real enough to require a control. The local names may change, but the boundary
between observed evidence, derived representation, hypothesis, and blocked
release should not disappear.

\section{External Review Positioning}
For external review, the paper should be positioned as a bounded methods result
rather than as a comprehensive account of Nwagu Aneke. A reviewer in digital
humanities should be able to evaluate whether the source/derived distinction is
documented and whether the article avoids overclaiming. A reviewer in
information systems should be able to evaluate whether the governance mechanism
is explicit enough to prevent claim drift. A reviewer in human-computer
interaction or design should be able to evaluate whether the proposed control
changes the behavior of generated interfaces or manuscripts. A reviewer in
African writing systems should be able to see what is claimed about the artifact
and what is explicitly not claimed.

This positioning matters because the paper is interdisciplinary but not
unbounded. It does not ask every reviewer to accept every adjacent field. It
offers each field a different inspection point. The artifact scholar inspects
the source boundary. The systems reviewer inspects the invariant or audit
method. The standards reviewer inspects interoperability. The ethics reviewer
inspects rights and authority boundaries. The AI reviewer inspects whether the
workflow reduces layer-promotion errors. A paper that can be attacked along
specific lines is stronger than one whose contribution dissolves into broad
language.

\section{Failure Cases That Remain Out of Scope}
Several important failures remain outside the approval. The paper does not
evaluate model performance on source images. It does not measure human reading
accuracy. It does not compare competing glyph-shape grammars. It does not
provide a public corpus. It does not settle dialect, orthographic, or
manuscript-historical questions. It does not claim that the derived layer is
better for teaching or software than the primary layer. It only states the layer
conditions under which such studies could be conducted.

Leaving these failures out of scope is not evasion. It is stage discipline. If a
future paper studies OCR, it needs image rights, crops, labels, baselines, and
error analysis. If a future paper studies tokenization, it needs corpora,
baselines, downstream tasks, and confidence intervals. If a future paper studies
Unicode readiness, it needs repertoire stability, glyph evidence, usage
examples, names, and community review. This paper provides a control condition
for those studies; it does not replace them.

\section{Approval Rationale}
The article is internally approvable because its claim is bounded, its evidence
is named, its limitations are active, and its approval does not require external
release. The central result can be checked against a finite set of files. The
manuscript avoids broad claims about universal compression, exceptional
mathematics, public source release, or completed standardization. The review
trace confirms that the paper is coherent as an internal research paper, while
the final readiness file keeps external submission separate.

This is the correct intermediate state for a serious lab. A weaker process would
either reject the paper because it is not externally complete or promote it as if
internal approval were public readiness. The stronger process records that the
paper is approved for the lab's research archive and still has additional work
before journal submission. That distinction allows the lab to accumulate real
papers without pretending that every internal result is already a venue-ready
article.

\section{Boundary Conditions for Future Promotion}
Future promotion from internal approval to external submission requires a
different gate. The first condition is venue-specific prior art: the paper must
be revised against the expectations of one target venue rather than remaining
generically interdisciplinary. The second condition is final source review:
claims about the artifact must be checked against the source dossier and any
available primary records. The third condition is rights and authority review:
the submission must identify exactly what source material, if any, becomes
public. The fourth condition is independent reviewer stress testing: at least one
reviewer should try to reject the paper on novelty, method, evidence, and ethics
grounds.

These conditions are not blockers to internal approval because the internal
paper does not claim external readiness. They are the route to responsible
publication. The lab can therefore keep researching from a stable foundation
without confusing internal scientific progress with journal acceptance. That
separation is central to the paper's method and to the broader Etisiobi research
standard.

\section{Practical Checklist for Reuse}
The practical checklist is short. A researcher reusing this paper should first
identify the layer of every claim. They should then list the transformation that
created any derived object. They should state the maximum public claim allowed
by that layer. They should add a negative control that would catch accidental
promotion. They should write limitations before writing contribution language.
They should record rights and authority assumptions separately from technical
results. Finally, they should rerun the approval check after every major
manuscript revision.

This checklist is deliberately operational. It translates the argument into a
repeatable workflow that other branches in the lab can use. If a future article
about OCR, tokenization, Unicode readiness, or design grammar cannot complete
this checklist, it should remain a research branch rather than an approved
paper. If it can complete the checklist, it becomes a candidate for the same
internal approval standard used here. The value of this paper is therefore not
only the local finding; it is the reusable discipline it imposes on future
research.

\section{Minimum Acceptance Claim}
The minimum acceptance claim is intentionally plain: the paper gives the lab a
checked way to keep an artifact-derived result useful without overstating its
source status. That is enough for internal approval, and it is also the claim
that future external reviewers should be asked to evaluate first.



\section{Standards Comparison}
The count-layer audit is compatible with existing standards rather than a
replacement for them. TEI can represent glyph declarations, uncertain readings,
and editorial apparatus; IIIF can identify canvases and regions; Web Annotation
can attach claims, comments, or classifications to source targets; PROV-O can
describe activities, agents, entities, and derivations; RO-Crate can package the
research object; DataCite can describe a citable release; CIDOC CRM can place
the artifact, interpretation event, and actor relationships inside a
cultural-heritage ontology; FAIR and CARE establish reuse and community
governance expectations \cite{azuonye1992,omniglotNwagu,unicodeAfricanScripts2023,scriptSourceNwagu,sirisNwaguProposal,teiP5Glyphs,iiifPresentation3,w3cAnnotationModel,w3cAnnotationVocab,w3cProvO,rocrate12,cidocCrm,datacite45,wilkinson2016fair,carroll2020care,wadden2020scifact,gao2023alce,karpathyAutoresearch,acmSubmissions,arxivSubmitTex}. None of these standards is
replaced by the audit. The audit supplies a decision that those standards can
carry: which count is primary-source-layer, which count is derived, and which
claims remain blocked.

This comparison matters because an underdocumented script can be harmed by
over-formalization. A premature data model may appear neutral while silently
choosing one interpretation as canonical. A premature package may make an
uncertain chart look like a complete dataset. A premature publication may cite a
primary source for a statement introduced by a later normalization step. The
audit is a corrective. It does not make the archive complete; it makes
incompleteness visible.

\section{Implications for Digital Humanities}
Digital humanities projects often turn source fragments into structured data.
That conversion is necessary for search, comparison, visualization, and
preservation, but it is also a site of epistemic risk. The Nwagu Aneke case
shows why source-critical boundaries must be represented as data rather than as
informal caveats. If the primary layer is 26 by 8 and the derived design layer
is 27/216, a database should not store only a single integer called "base
count." It should store the evidence layer, transformation rule, review status,
and allowed claim.

The practical consequence is that every downstream system receives a safer
input. A digital edition can expose the 26 by 8 grid without suppressing derived
analytical forms. A teaching interface can use the f/v split while identifying
it as a derived teaching layer. A standards-readiness matrix can ask which layer
is stable enough for encoding discussion. A computational benchmark can test
whether models preserve the distinction. This is the sense in which the
count-layer audit becomes infrastructure: it is not merely a finding about one
number, but a constraint on later representations.

\section{Reviewer Objections and Responses}
A skeptical reviewer may object that the result is too small. The response is
that exact-count drift is not small when downstream theory, software, and public
claims depend on the count. In a mature archive, preventing one false foundation
claim can be more valuable than adding another speculative mapping. The paper's
claim is deliberately narrow because the evidence demands it.

A second objection is that the audit depends on the current evidence package.
That is correct. The paper is not a final critical edition. It is a reproducible
audit of the present layer distinction. Its strength is that it specifies how it
can be overturned. If a reviewed source changes the primary layer, the ledger
and paper must change. A result with clear revision conditions is stronger than
a larger theory with no falsification path.

A third objection is that many standards already handle provenance and
annotation. That is also correct. The article does not claim to invent
provenance or annotation. It contributes an artifact-specific decision rule that
works with those standards. Standards can encode an assertion; they do not
automatically know whether the assertion over-promotes a derived layer.

\section{Applications Enabled by the Audit}
The audit enables several applications without approving them prematurely. The
first is a source-critical digital edition in which every row, column, cell, and
derived split is represented with a layer label. The second is an OCR or
annotation workflow in which model outputs cannot upgrade uncertain cells into
source facts. The third is a tokenizer or design grammar that may use the
derived matrix as a normalization while preserving the primary ledger. The
fourth is an article-selection system that blocks manuscripts whose abstracts
collapse the layers.

These applications are not claimed as completed results in this paper. They are
listed because the audit makes them testable. For each application, the question
is no longer "can we use Nwagu Aneke to inspire a system?" The question becomes
"which layer does the system use, what claim ceiling follows from that layer,
and what negative control would catch an over-promotion error?" That change in
question is the route from symbolic inspiration to research-grade systems.

\section{Ethical and Community Boundary}
The paper treats the artifact as culturally meaningful source material, not as a
free pool of symbols for extraction. CARE principles are therefore relevant even
when no public image release is attempted \cite{carroll2020care}. The text
does not reproduce source images, publish a manuscript corpus, assign public
codepoints, or claim community authorization for a release. It makes a bounded
claim about count-layer discipline inside the current research package.

This boundary is part of the method. A weaker paper would hide the cultural and
rights boundary as a limitation paragraph after making broad claims. This paper
keeps the boundary active throughout the argument. If an application requires
public source images, manuscript access, community review, or an encoding
proposal, it must pass a separate gate. The present approval is internal and
textual.


\section{Article-Specific Implementation Detail}
The implementation detail differs by article, but the approval rule is the same.
The article-specific result must be represented as a bounded object with a named
experiment, a claim ceiling, and an explicit non-claim. For this article, the
result is not treated as evidence for a larger theory. It is treated as a
controlled contribution to provenance, Indigenous data governance, and responsible cultural-heritage release. The public-facing experiment label is the provenance-ledger run, and the manuscript is approved only to the extent that it preserves the result
statement and the source/derived foundation.

This matters for the remaining Nwagu Aneke portfolio. The lab can maintain many
branches, but each branch must earn approval through a specific result. A source
ledger is not the same as a tokenizer baseline. A selector mapping is not the
same as a Unicode proposal. A provenance ledger is not the same as public corpus
release. A comparative matrix is not the same as standardization. Each article
must therefore make one contribution and refuse the adjacent stronger claims.

The implementation detail also makes future revision cheaper. If later review
changes the source foundation, the article's approval can be rechecked. If later
experiments strengthen the result, the paper can move toward external
submission. If later rights review blocks public discussion, the approval can be
withdrawn without rewriting the whole portfolio. This is what internal approval
is for: a stable, inspectable state between working notes and public submission.

\section{Internal Approval Value}
The internal approval value is practical. The paper can be cited inside the lab,
used as a basis for the next experiment, and compared against later revisions
without pretending that a journal has accepted it. It gives collaborators one
stable object to inspect: the claim, the evidence, the method, the limits, and
the next test. That stability is what lets the research program compound.

The paper also becomes a unit of comparison. Later runs can ask whether a new
experiment improves the claim, narrows the limitation, changes the rights
boundary, or overturns the result. If none of those changes occurs, the article
should not be rewritten merely to sound stronger.

That comparison rule is deliberately conservative. It prevents the portfolio
from treating parallel articles as interchangeable proof of one master theory.
Each article must carry its own local burden: a specific artifact relationship,
a specific experiment output, a specific source ceiling, and a specific
publication risk. When the burden is local, approval can be audited locally; when
the burden is vague, the article is pushed back to the research-program stage.

This is also a coordination device for the lab. The same source object can
support several research directions, but each direction needs a different
acceptance test. By forcing the acceptance test into the article, the paper tells
future researchers what would improve it, what would invalidate it, and what
must remain outside the claim. That is the difference between a research article
and a portfolio note.

\section{Limitations}
source transcription review remains required before any stronger source-layer statement, and rights review remains required before any public source-image, corpus, or authority claim. 

This paper is internally approved only within its bounded scope. It does not
authorize public source-image release. It does not claim a completed Unicode
proposal, full glyph-shape grammar, manuscript corpus, or public cultural
authority decision. It preserves the accepted source foundation and treats the
derived f/v layer as derived.

\section{Reproducibility}
The analysis package records the manuscript source, evidence ledger, article-specific result, review notes, rights boundary, and reproducibility materials needed to inspect the claim without treating the paper as a public release. 

The analysis package includes the manuscript source, experiment result, claim
audit, novelty audit, source-review gate, rights-authority gate,
reproducibility packet, and review-team trace. Reproduction means checking the
experiment result, checking that the manuscript preserves the 26 by 8 and
derived f/v distinction, and checking that no public-release or broad-theory
claim has been inserted.

\section{Conclusion}
This article contributes a bounded internal research result for provenance, Indigenous data governance, and responsible cultural-heritage release. It is
approved as an Etisiobi research paper because it has a result, a method, a
claim ceiling, review files, and reproducibility materials. It is not an
external submission package. Its value is that it extends the Nwagu Aneke
research program without promoting derived, speculative, or blocked claims into
primary source claims.

\bibliographystyle{ACM-Reference-Format}
\bibliography{../references}
\end{document}



